\chapter{Introduction}
Nonadiabatic dynamics simulations are essential for understanding a wide range of chemical processes, including photochemical reactions, energy transfer, and charge transfer events. In traditional theoretical chemistry, the behaviour of molecular systems is often described using the \acrfull{boa}, which assumes that the motion of nuclei and electrons can be separated due to their vastly different masses. However, in many cases, this approximation breaks down, particularly when electronic states are closely spaced. In these situations, nonadiabatic effects become significant, and the dynamics of the system cannot be accurately described without considering the coupling between electronic and nuclear degrees of freedom.\supercite{acs.chemrev.7b00577}

While the exact solution of the \acrfull{tdse} for nonadiabatic systems provides a complete description of the dynamics, applying this approach to realistic molecular systems is computationally prohibitive due to the exponential scaling of the required computational resources with system size. To overcome this challenge, chemists usually rely on a variety of approximate methods that aim to capture the essential physics of nonadiabatic processes while remaining computationally feasible. These methods include mixed quantum-classical dynamics approaches, such as \acrfull{tsh}, which treat the nuclei classically while maintaining a quantum mechanical description of the electrons. Additionally, various numerically exact methods have been developed to provide benchmark results for small systems, allowing for the assessment of the accuracy of approximate techniques.

This thesis presents a comprehensive study of computational strategies for simulating nonadiabatic dynamics, with a focus on the development, application and validation of novel algorithms. By combining rigorous mathematical theory with complex software engineering, I aim to provide a deep understanding of the computational tools required to simulate nonadiabatic processes where \acrshort{boa} breaks down.

The theoretical foundations of the methods discussed in this thesis are derived in detail to ensure a clear understanding of the underlying principles and assumptions. This includes the strict mathematical framework of Hilbert spaces, exact numerical methods for solving the \acrfull{tdse} like the Split-Operator method or Crank--Nicolson scheme, and the theory of mixed quantum-classical dynamics with a particular focus on \acrshort{tsh}. Crucially, all of the exact quantum dynamics methods and techniques described in these theoretical sections were implemented by me in a custom software package \textsc{Zinq}.\supercite{10.5281/zenodo.18386143}

Building on these theoretical foundations, I apply the developed methods to a range of model systems and real molecular systems to investigate their nonadiabatic behaviour. Chronologically, the first primary scientific contribution of this thesis is a critical evaluation of the extreme sensitivity of \acrshort{tsh} to underlying electronic structure methods. By investigating the performance of \acrshort{tsh} across a variety of electronic structure methods, I demonstrate that the choice of method can significantly impact the predicted nonadiabatic dynamics, leading to qualitatively different outcomes. This highlights the importance of carefully selecting and validating electronic structure methods when using \acrshort{tsh} for nonadiabatic simulations.

The second major contribution of this thesis investigates curvature-driven surface hopping algorithms, which attempt to bypass the need for explicitly calculated nonadiabatic couplings. Applying these algorithms to a range of model and molecular systems, I assess their performance in capturing the accurate \acrfull{tdc} and the resulting dynamics. The results show that while curvature-driven algorithms can provide a computationally efficient alternative to traditional \acrshort{tsh}, they may not always capture the correct physics of nonadiabatic processes, particularly in systems with complex electronic structure.

Finally, the scope of the research is broadened to atomic collision dynamics governed by strong spin-orbit coupling, where multidimensional wavepacket propagation and analytical transition theories are integrated to model complex charge transfer events. On the high-level \acrfull{pes} of the collision system, I apply the exact solution to \acrshort{tdse} to obtain scattering probabilities and cross sections, which are then compared to results obtained from analytical transition theories. This work provides a novel approach to modeling atomic collision dynamics in the presence of strong spin-orbit coupling, offering insights into the underlying physical processes and the validity of different theoretical approaches.
